\documentclass[UTF8]{ctexart}
\usepackage[a4paper,margin=2.6cm]{geometry}
\begin{document}

\section*{第六部分：距离顺序问题的下界}

这一节证明距离顺序问题的下界。作者先证明一个时间下界，再证明两个比较次数下界，最后把它们合并起来。

\subsection*{6.1 时间下界}

Lemma 6.1 说明：在时间模型中，任何正确的确定性算法，在任意图上解决距离顺序问题都需要 \(\Omega(m)\) 时间。

证明使用的是反证和对抗构造。设 \(L=[v_1,v_2,\ldots,v_n]\) 是图 \(G\) 的某个距离顺序。作者给弧设置长度，使得这个顺序 \(L\) 成为唯一的真实距离顺序。具体思想是：如果弧从 \(v_i\) 指向 \(v_j\)，就让它的长度和 \(j-i\) 的关系匹配，使得 \(v_i\) 的真实距离按下标递增。

然后考虑一个算法如果没有充分读取邻接表，会发生什么。如果算法完全没有访问某个较早顶点的出弧列表，那么可以偷偷加入或缩短一条从这个顶点跳到更后顶点的弧，使得真实距离顺序发生改变，但算法由于没有看过这部分输入，行为仍然和原图上一样，于是会输出错误顺序。

类似地，如果算法虽然访问了某个邻接表，但没有访问完，也可以把未访问的最后一条弧改短，制造一个新的图，让正确的距离顺序改变，而算法仍然无法察觉。

因此，正确算法必须花足够多的时间去查看输入。结合 \(n>2\) 的假设，可以得到 \(\Omega(m)\) 时间下界。

对无向图，证明略有不同，但结论类似。因为每条无向边对应两个方向的弧，算法至少要检查足够多的边，否则同样可能漏掉一条会改变距离顺序的前向弧。

\subsection*{6.2 比较次数下界之一：\(\log D\)}

Lemma 6.2 给出第一个比较下界：在比较模型中，任何正确的确定性算法，在最坏情况下至少需要 \(\lceil \log D\rceil\) 次比较；在某个最坏输入分布下，期望比较次数至少为 \(\lfloor \log D\rfloor\)。

这里的 \(D\) 是图的距离顺序数量。证明思路和排序下界类似：一个算法的比较过程可以看成决策树。每一种可能的距离顺序都需要被算法区分出来，因此决策树至少要有 \(D\) 个叶子。二叉决策树若要区分 \(D\) 种情况，高度至少是 \(\log D\)。

论文中还处理了比较结果等于 0 的情况。作者通过对弧长做任意小的扰动，保证真实距离顺序不变，同时避免比较结果正好为 0。这样就能把问题归约到标准的二分决策树下界。

\subsection*{比较次数下界之二：\(F-n+1\)}

Lemma 6.4 给出第二个比较下界：在比较模型中，任何正确的确定性算法至少需要
\[
F-n+1
\]
次比较。

这里 \(F\) 是某个距离顺序中最多可能出现的前向弧数量。证明的核心思想是：如果算法比较太少，就无法充分判断哪些前向弧会真正影响真实距离顺序。

作者选取一个前向弧数为 \(F\) 的距离顺序 \(L=[v_1,v_2,\ldots,v_n]\)，并给前向弧设置长度，使 \(L\) 成为唯一真实距离顺序。对非前向弧，则统一设置成很大的长度，使它们不会影响算法行为。接着，把算法看成一棵只依赖前向弧长度的决策树。

如果算法在某条路径上只做了不超过 \(F-n\) 次比较，那么这些比较只给出了不超过 \(F-n\) 个线性约束。再加上相邻前向弧的约束，总约束数仍少于前向弧变量数，因此仍然存在自由度。沿着这个自由度改变弧长，可以让某条非相邻前向弧的长度变为 0，同时保持算法看到的所有比较结果不变。

这样构造出的新弧长会改变真实距离顺序：某个后面的顶点会因为零长度跳边而必须提前出现。但算法的比较结果完全一样，所以它仍然输出原来的 \(L\)，从而出错。于是，正确算法必须至少做 \(F-n+1\) 次比较。

对无向图，由于每条边在任意顺序下恰好对应一条前向弧，所以 \(F=m\)，下界变成 \(m-n+1\) 次比较。

\subsection*{合并下界}

时间下界和比较下界可以合并。因为比较本身也需要时间，所以 Lemma 6.1 和 Lemma 6.2 得到：

\[
\Omega(m+\log D)
\]

这是时间模型中的下界。

比较模型里，Lemma 6.2 和 Lemma 6.4 合并得到：

\[
\Omega(F-n+1+\log D)
\]

这是有向图上的比较次数下界。对无向图，因为 \(F=m\)，对应下界为

\[
\Omega(m-n+1+\log D).
\]

\subsection*{这一节的意义}

第 6 节的作用是给后面证明 Dijkstra 最优性打基础。作者不是只证明 Dijkstra 很快，而是先证明任何算法都不可能比这些下界快太多。后面如果 Dijkstra 能达到同阶上界，就能说明它在相应模型中是 universally optimal 的。

\end{document}
